Quantile trend filtering estimates time-varying conditional quantiles under piecewise polynomial structure, and its state-of-the-art frequentist solution is an ADMM algorithm applied to the check loss with an $\ell_1$ difference penalty \citep{brantley2020baseline}. That procedure is an optimizer, not a statistical sampler: it returns a single point estimate at a prespecified penalty level, with no uncertainty quantification and no principled mechanism for selecting the regularization parameter or the likelihood scale. This paper develops the fully Bayesian counterpart of that estimator. Starting from the model whose posterior mode is the ADMM solution---an asymmetric Laplace likelihood paired with a global Laplace prior on the trend differences---we build a hierarchy in which the regularization parameter and the likelihood scale receive conjugate priors and are inferred jointly with the trend. The resulting three-block Gibbs sampler updates both scale parameters in closed form and updates the trend via a single Proximal Metropolis-Adjusted Langevin (P-MALA) step \citep{pereyra2016proximal}, whose proximal map is itself a trend filtering problem solved by a small number of ADMM iterations. Because the inner ADMM coefficient matrix is fixed, sparse, and banded, its factorization is computed once and reused, giving per-iteration cost linear in the sample size. The optimization machinery of the frequentist estimator is thus repurposed as the computational engine of a posterior sampler.

The framework extends to joint estimation of multiple quantiles under exact non-crossing constraints, enforced through box projections embedded in the ADMM proximal operator, with each quantile updated sequentially by its own P-MALA step. For extreme probability levels, where the near-constant gradient of the asymmetric Laplace likelihood causes gradient-based samplers to drift, we introduce a sentinel quantile strategy: auxiliary quantiles placed beyond the target range whose box constraints confine the target within a dynamically updating interval, preventing drift while preserving exploration.

Simulation studies spanning piecewise constant, varying smoothness, and smooth signals under Gaussian, Beta, and mixed normal noise show that the proposed method delivers accurate, non-crossing quantile estimates across interior and extreme probability levels, with full posterior uncertainty quantification obtained as a direct byproduct of inference. Compared with the strongest existing Bayesian sampler for this problem \citep{onizuka2024fast}, whose per-iteration cost is cubic in the sample size and whose accuracy degrades at extreme quantiles, the proposed sampler is competitive or superior in estimation accuracy, scales linearly, and resolves the extreme-quantile instability.
